% ──────────────────────────────────────────────────
%  part1.tex：內容範例（示範版面樣式，可整段替換）
% ──────────────────────────────────────────────────
\section{段落標題}
\subsection{小節標題}

\begin{narrowframed}
\textbf{定義 1.1}\textbf{ 名稱}\\
內容文字，數學式如 $x^2 + y^2 = r^2$。
\end{narrowframed}

\begin{itemize}
    \item 項目一
    \item 項目二
\end{itemize}

\begin{whitebox}
無框區塊，用來隔開段落內容。
\end{whitebox}

\subsection{表格範例}

\begin{table}[H]
    \centering
    \caption{一般表格}
    \vspace{8pt}
    \begin{tabular}{lcc}
        \toprule
        欄位 & 數值 A & 數值 B \\
        \midrule
        項目一 & 1.00 & 2.00 \\
        項目二 & 3.00 & 4.00 \\
        \bottomrule
    \end{tabular}
\end{table}

\begin{table}[H]
    \centering
    \caption{寬表格，以 resizebox 縮放至版寬}
    \vspace{8pt}
    \resizebox{\textwidth}{!}{%
    \begin{tabular}{lcccccc}
        \toprule
        名稱 & 第一欄位 & 第二欄位 & 第三欄位 & 第四欄位 & 第五欄位 & 第六欄位 \\
        \midrule
        很長的項目名稱 & 1.000 & 2.000 & 3.000 & 4.000 & 5.000 & 6.000 \\
        \bottomrule
    \end{tabular}%
    }
\end{table}

\subsection{習題範例}

\begin{examplebox}{1}{題目名稱}
題目敘述。
\end{examplebox}

\begin{solutionbox}
解題過程。
\end{solutionbox}

\begin{analysisbox}
解題思路與考點。
\end{analysisbox}

\subsection{程式碼與多欄範例}

\begin{lstlisting}[language=R]
lm_model <- lm(y ~ x, data = df)
summary(lm_model)
\end{lstlisting}

\noindent
\begin{minipage}[t]{0.48\textwidth}
左欄內容。
\end{minipage}
\hfill
\begin{minipage}[t]{0.48\textwidth}
右欄內容。
\end{minipage}

% \begin{figure}[H]
%     \centering
%     \includegraphics[width=0.7\textwidth]
%         {content/__SUBJECT__/note/note_1/figures/檔名}
%     \caption{圖片標題}
% \end{figure}
